% \iffalse meta-comment
%
% hit-listing.dtx 是各个类共用的代码清单层
%
% \fi
%
% \section{代码清单}
%
% \changes{v3.2a}{2026/10/06}{代码清单改由文档类排：\pkg{listings} 在导言区末尾由类装上，写了题注的
%   lstlisting 按 raw-style 设置外框与行号（\option{raw-style} 的
%   \option{frame} 与 \option{line-numbering}），题注照表题排、名字“代码”，caption 接受主副语言内容，
%   续排不占新号、行号接着数（\texttt{continued}）；没写题注的是正文里的代码块，
%   不加框、不编行号。代码用等宽字体、正文字号，行距照正文段的倍数按等宽字体的行高算}
%
% 有题注的 lstlisting 按代码清单排版，无题注时作为普通代码块。
% \begin{itemize}
% \item \option{frame}：\texttt{tabular}（默认）题注在上、居中，下面粗线、代码、粗线；
%   \texttt{ruled} 粗线、居左的题注、细线、代码、粗线；\texttt{boxed} 一圈细线的框，题注在框下居中；
%   \texttt{none} 不画线，题注在上。粗细取三线表那两个（\option{table-rules}）。
% \item \option{line-numbering}：行号的样子，写法与伪代码的同一个键一样（首字定数字的写法，
%   后面跟着的照印）；\texttt{none} 不编行号。默认 \texttt{1}。
% \end{itemize}
%
% 代码使用等宽字体和正文字号；行号撑杆须与代码行盒一致，避免额外撑高。
%
% \pkg{listings} 的续页不会自己标“（续代码）”，跨页的清单照原样断开；要分成两张的，后一张写
% \texttt{continued}。文档自己装了 \pkg{minted} 这类别的包，\env{lstlisting} 以外的环境类不管。
%    \begin{macrocode}
%<*shared-listing>
%    \end{macrocode}
%
% \subsection{键}
%
% 两个键的值各存一份全局 \texttt{tl}，每开一个清单现读。\texttt{auto} 退回默认值。
%    \begin{macrocode}
\tl_new:N \g__hit_listing_frame_tl
\tl_new:N \g__hit_listing_numbering_tl
\tl_gset:Nn \g__hit_listing_frame_tl { tabular }
\tl_gset:Nn \g__hit_listing_numbering_tl { 1 }
\keys_define:nn { hit / style }
  { listing .code:n = { \hit@keys@set:nn { hit / style / listing } {#1} } }
\keys_define:nn { hit / style / listing }
  {
    frame .choices:nn = { tabular , ruled , boxed , none , auto }
      {
        \str_if_eq:VnTF \l_keys_choice_tl { auto }
          { \tl_gset:Nn \g__hit_listing_frame_tl { tabular } }
          { \tl_gset_eq:NN \g__hit_listing_frame_tl \l_keys_choice_tl }
      } ,
    frame .value_required:n = true ,
    line-numbering .code:n =
      {
        \str_if_eq:nnTF {#1} { auto }
          { \tl_gset:Nn \g__hit_listing_numbering_tl { 1 } }
          { \tl_gset:Nn \g__hit_listing_numbering_tl {#1} }
      } ,
    line-numbering .value_required:n = true ,
  }
%    \end{macrocode}
%
% \subsection{装包}
%
% 导言区末尾没装 \pkg{listings} 就装上。设定挂在它的装后钩子上，文档自己在导言区装了
% 再写 \cs{lstset}，写的照样压过这里的。\hologo{pdfLaTeX} 下中文是几个字节的活动字符，
% \texttt{extendedchars} 会把它们当成单个 8 位字符，报 Improper alphabetic constant，所以关掉，
% 原样交给 \pkg{CJK} 排；定宽的列也要逐字节取字符码，同样报错，那边改成 \texttt{flexible}（\texttt{fullflexible} 会丢行首缩进）。
% 字符串里的空格照排空格，不印 \pkg{listings} 默认的“␣”。
%    \begin{macrocode}
\AddToHook { begindocument / before }
  { \@ifpackageloaded { listings } { } { \RequirePackage { listings } } }
\AddToHook { package / listings / after } { \__hit_listing_package: }
\cs_new_protected:Npn \__hit_listing_package:
  {
    \lstset
      {
        basicstyle = \__hit_listing_font: ,
        keywordstyle = { } ,
        columns = [l] fixed ,
        basewidth = { \__hit_listing_basewidth: , 0.45em } ,
        keepspaces = true ,
        showstringspaces = false ,
        breaklines = true ,
        aboveskip = 0pt ,
        belowskip = 0pt ,
        abovecaptionskip = 0pt ,
        belowcaptionskip = 0pt ,
        numbers = left ,
        numberstyle = \__hit_listing_number_font: ,
        numberfirstline = true ,
        everydisplay = \__hit_listing_display: ,
      }
    \sys_if_engine_pdftex:T { \lstset { extendedchars = false , columns = flexible } }
    \lst@Key { caption } \relax { \hit@bilingual@listing@caption:n {##1} }
    \lst@Key { continued } { false } [ true ]
      {
        \str_if_eq:nnTF {##1} { true }
          { \bool_set_true:N \l__hit_listing_continued_bool }
          { \bool_set_false:N \l__hit_listing_continued_bool }
      }
    \lst@AddToHook { PreSet } { \__hit_listing_preset: }
    \lst@AddToHook { PreInit } { \__hit_listing_preinit: }
    \lst@AddToHook { EveryPar } { \__hit_listing_strut: }
    \lst@AddToHook { DeInit } { \__hit_listing_deinit: }
    \AddToHook { begindocument / end } { \__hit_listing_counter: }
  }
\tl_new:N \l__hit_listing_en_tl
\bool_new:N \l__hit_listing_continued_bool
%    \end{macrocode}
%
% 编号照图表：学位论文按章（“1-1”），报告不排章，从 1 数下去。\pkg{listings} 在
% \cs{AtBeginDocument} 里才建计数器，所以挪到 \texttt{begindocument/end}。
%
% 题注也在那时接管：\pkg{caption}（由 \pkg{subcaption} 带进来）在 \cs{AtBeginDocument} 里把
% \pkg{listings} 的 \cs{lst@MakeCaption} 换成自己的一套，题注名取 \cs{lstlistingname}，排法走
% \cs{@makecaption}。这里换回 \pkg{listings} 原来那个（\pkg{caption} 存着），题注交给本节的
% \cs{\_\_hit\_listing\_makecaption:nn}。
%    \begin{macrocode}
\cs_new_protected:Npn \__hit_listing_counter:
  {
    \cs_if_exist:NT \caption@ORI@lst@MakeCaption
      { \cs_set_eq:NN \lst@MakeCaption \caption@ORI@lst@MakeCaption }
    \cs_set_eq:NN \lst@makecaption \__hit_listing_makecaption:nn
    \bool_if:NT \g__hit_chapter_bool
      { \counterwithin* { lstlisting } { chapter } }
    \cs_set:Npn \thelstlisting
      {
        \bool_if:NT \g__hit_chapter_bool { \thechapter - }
        \int_use:N \c@lstlisting
      }
    \hit@numbering@counter@init:n { lstlisting }
  }
%    \end{macrocode}
%
% \subsection{字体与行距}
%
% 代码块沿用等宽字宽校正，行盒按目标字体在正文字号上的度量计算。
%    \begin{macrocode}
\dim_new:N \l__hit_listing_size_dim
\dim_new:N \l__hit_listing_lead_dim
\tl_new:N \l__hit_listing_natural_tl
\tl_new:N \l__hit_listing_ascent_tl
\dim_new:N \l__hit_listing_ascent_dim
\cs_new_protected:Npn \__hit_listing_metrics:
  {
    \dim_set_eq:NN \l__hit_listing_size_dim \g_docxlike_body_size_dim
    \dim_compare:nNnT { \l__hit_listing_size_dim } = { 0pt }
      { \dim_set:Nn \l__hit_listing_size_dim { \f@size pt } }
    \tl_set_eq:NN \l__hit_listing_natural_tl \g_docxlike_mono_single_tl
    \tl_set_eq:NN \l__hit_listing_ascent_tl \g_docxlike_mono_ascent_tl
    \dim_set:Nn \l__hit_listing_ascent_dim
      { \l__hit_listing_ascent_tl \l__hit_listing_size_dim + \g_docxlike_mono_ascent_offset_dim }
    \group_begin:
      \docxlike_format_compute:n { Normal }
      \str_if_eq:VnTF \l_docxlike_format_line_unit_tl { exact }
        {
          \dim_gset:Nn \g_tmpa_dim { \l_docxlike_format_line_value_tl }
          \dim_gset:Nn \g_tmpb_dim { .8 \g_tmpa_dim }
        }
        {
          \dim_gset:Nn \g_tmpa_dim
            {
              \fp_to_dim:n
                {
                  ( \str_if_eq:VnTF \l_docxlike_format_line_unit_tl { auto }
                      { \l_docxlike_format_line_value_tl } { 1 } )
                  * \l__hit_listing_natural_tl * \l__hit_listing_size_dim
                }
            }
          \dim_gset_eq:NN \g_tmpb_dim \l__hit_listing_ascent_dim
        }
    \group_end:
    \dim_set_eq:NN \l__hit_listing_lead_dim \g_tmpa_dim
    \dim_set_eq:NN \l__hit_listing_ascent_dim \g_tmpb_dim
  }
\cs_generate_variant:Nn \docxlike_font_metric:nNNF { V }
%    \end{macrocode}
%
% 字宽为等宽字体步进加半份字符网格增量，每个字符在格内左对齐。
%    \begin{macrocode}
\cs_new:Npn \__hit_listing_basewidth:
  {
    \tex_dimexpr:D \tex_fontcharwd:D \tex_font:D `0
      \dim_compare:nNnT { \g_docxlike_page_char_dim } > { 0pt }
        { + ( \g_docxlike_page_char_dim - \g_docxlike_body_size_dim ) / 2 }
    \scan_stop:
  }
\cs_new_protected:Npn \__hit_listing_font:
  {
    \__hit_listing_metrics:
    \fontsize
      { \l__hit_listing_size_dim }
      { \l__hit_listing_lead_dim }
    \ttfamily \selectfont
    \skip_zero:N \lineskip
    \dim_zero:N \lineskiplimit
    \hbox_set:Nn \strutbox { \__hit_listing_strut: }
  }
\cs_new_protected:Npn \__hit_listing_number_font:
  {
    \fontsize
      { \fp_to_dim:n { 0.8 * \l__hit_listing_size_dim } }
      { \l__hit_listing_lead_dim }
    \ttfamily \mdseries \upshape \selectfont
  }
%    \end{macrocode}
%
% \subsection{每一张}
%
% \texttt{PreSet} 在读选项之前：清掉上一张的英文题与续排标志，外框定题注的位置。
% \texttt{everydisplay} 在读完选项之后、题注之前：写了题注（或 \texttt{continued}）的是清单，
% 照外框缩进单元格边距、上面排空段；续排的短题注清空，\pkg{listings} 见了就不步进计数器、
% 不进清单，\cs{label} 拿到的仍是上一张的号，行号接着上一张。代码的每一行在 \pkg{listings} 里都是
% 一段，\pkg{docxlike} 的段落钩子会给它发正文的撑杆（宋体的行盒比等宽字体的高），所以把那两个钩子
% 关掉，用的是它给页眉页脚留的开关 \cs{l\_docxlike\_format\_in\_header\_bool}，出清单的组即还原。\texttt{PreInit} 在题注之后、头一行之前：没写题注的是正文代码块，不编行号。
% \texttt{EveryPar} 每行开头放一根撑杆，行盒的高深就是等宽字体的上升部与行距减上升部，
% 一行接一行、贴着线排。
%    \begin{macrocode}
\bool_new:N \l__hit_listing_captioned_bool
\dim_new:N \l__hit_listing_width_dim
\cs_new_protected:Npn \__hit_listing_preset:
  {
    \tl_clear:N \l__hit_listing_en_tl
    \bool_set_false:N \l__hit_listing_continued_bool
    \str_if_eq:VnTF \g__hit_listing_frame_tl { boxed }
      { \def \lst@captionpos { b } }
      { \def \lst@captionpos { t } }
  }
\cs_new_protected:Npn \__hit_listing_display:
  {
    \dim_set_eq:NN \l__hit_listing_width_dim \linewidth
    \bool_set_true:N \l_docxlike_format_in_header_bool
    \__hit_listing_metrics:
    \bool_if:NT \l__hit_listing_continued_bool
      {
        \tl_if_empty:NT \lst@caption { \def \lst@caption { \scan_stop: } }
        \cs_set_eq:NN \lst@@caption \c_empty_tl
        \let \lst@firstnumber \c@lstnumber
      }
    \bool_set:Nn \l__hit_listing_captioned_bool
      { ! \tl_if_empty_p:N \lst@caption }
    \bool_if:NT \l__hit_listing_captioned_bool
      {
        \str_if_eq:VnF \g__hit_listing_frame_tl { none }
          {
            \edef \lst@xleftmargin { \dim_use:N \tabcolsep }
            \edef \lst@xrightmargin { \dim_use:N \tabcolsep }
          }
        \__hit_float_blank:n { table-above }
      }
  }
\cs_new_protected:Npn \__hit_listing_preinit:
  {
    \bool_lazy_and:nnTF
      { \l__hit_listing_captioned_bool }
      { ! \str_if_eq_p:Vn \g__hit_listing_numbering_tl { none } }
      {
        \edef \lst@numbersep { \dim_eval:n { 0.5 \l__hit_listing_size_dim } }
        \__hit_listing_number_style:
      }
      { \cs_set_eq:NN \lst@PlaceNumber \prg_do_nothing: }
    \bool_lazy_and:nnT
      { \l__hit_listing_captioned_bool }
      { \str_if_eq_p:Vn \g__hit_listing_frame_tl { boxed } }
      { \__hit_listing_box_begin: }
  }
\cs_new_protected:Npn \__hit_listing_strut:
  {
    \tex_vrule:D width 0pt
      height \l__hit_listing_ascent_dim
      depth \fp_to_dim:n
        { \l__hit_listing_lead_dim - \l__hit_listing_ascent_dim }
      \scan_stop:
  }
\cs_new_protected:Npn \__hit_listing_deinit:
  {
    \bool_if:NT \l__hit_listing_captioned_bool
      {
        \str_case:Vn \g__hit_listing_frame_tl
          {
            { tabular } { \__hit_listing_rule:n { \g__hit_table_rule_outer_tl } }
            { ruled } { \__hit_listing_rule:n { \g__hit_table_rule_outer_tl } }
            { boxed } { \__hit_listing_box_end: }
          }
        \str_if_eq:VnF \g__hit_listing_frame_tl { boxed }
          { \__hit_float_blank:n { table-below } }
      }
  }
\cs_new_protected:Npn \__hit_listing_rule:n #1
  {
    \par \nointerlineskip
    \tex_hrule:D width \l__hit_listing_width_dim height \dim_eval:n {#1} depth 0pt \scan_stop:
    \nointerlineskip
  }
%    \end{macrocode}
%
% 行号：写法照伪代码的 \option{line-numbering}（\cs{\_\_hit\_algorithm\_number:nn}）。
%    \begin{macrocode}
\tl_new:N \l__hit_listing_nl_style_tl
\tl_new:N \l__hit_listing_nl_suffix_tl
\cs_new_protected:Npn \__hit_listing_number_style:
  {
    \tl_set:Nn \l__hit_listing_nl_style_tl { 1 }
    \tl_set_eq:NN \l__hit_listing_nl_suffix_tl \g__hit_listing_numbering_tl
    \clist_map_inline:nn { ① , 一 , 1 , i , I , a , A }
      {
        \str_if_eq:eeT
          { \str_range:Nnn \g__hit_listing_numbering_tl { 1 } { \str_count:n {##1} } }
          { \tl_to_str:n {##1} }
          {
            \tl_set:Nn \l__hit_listing_nl_style_tl {##1}
            \tl_set:Ne \l__hit_listing_nl_suffix_tl
              { \str_range:Nnn \g__hit_listing_numbering_tl { \str_count:n {##1} + 1 } { -1 } }
            \clist_map_break:
          }
      }
    \cs_set:Npn \thelstnumber
      {
        \__hit_algorithm_number:VV \l__hit_listing_nl_style_tl { \int_use:N \c@lstnumber }
        \l__hit_listing_nl_suffix_tl
      }
  }
%    \end{macrocode}
%
% 框图型：一圈细线。整张装进一个竖盒量高，再画四边；框里上下各让出半个线宽，
% 左右是单元格边距。
%    \begin{macrocode}
\box_new:N \l__hit_listing_box
\cs_new_protected:Npn \__hit_listing_box_begin:
  {
    \par
    \setbox \l__hit_listing_box \vbox \bgroup
      \dim_set:Nn \hsize { \l__hit_listing_width_dim - 2 \dim_eval:n { \g__hit_table_rule_inner_tl } }
  }
\cs_new_protected:Npn \__hit_listing_box_end:
  {
    \par
    \egroup
    \dim_set:Nn \l_tmpa_dim { \g__hit_table_rule_inner_tl }
    \nointerlineskip
    \vbox:n
      {
        \tex_hrule:D height \l_tmpa_dim depth 0pt \scan_stop:
        \hbox:n
          {
            \tex_vrule:D width \l_tmpa_dim \scan_stop:
            \vbox:n
              {
                \skip_vertical:n { \l_tmpa_dim / 2 }
                \hbox_to_wd:nn { \l__hit_listing_width_dim - 2 \l_tmpa_dim }
                  { \box_use_drop:N \l__hit_listing_box \tex_hss:D }
                \skip_vertical:n { \l_tmpa_dim / 2 }
              }
            \tex_vrule:D width \l_tmpa_dim \scan_stop:
          }
        \tex_hrule:D height \l_tmpa_dim depth 0pt \scan_stop:
      }
    \nointerlineskip
  }
%    \end{macrocode}
%
% \subsection{题注}
%
% 题注照表题排（样式 \option{caption-table}），名字“代码”，题号与题名之间的间隔照图表。线条型的
% 题注照单元格的样式居左，前后是粗线与细线；表格型的题注下面接粗线。博士档写了英文题的，下面
% 再排一行“Listing 1-1”加英文题。续排的只排“代码1-1（续代码）”，居右，不占新号。
%    \begin{macrocode}
\cs_new_protected:Npn \__hit_listing_makecaption:nn #1#2
  {
    \def \@captype { lstlisting }
    \str_if_eq:VnT \g__hit_listing_frame_tl { ruled }
      { \__hit_listing_rule:n { \g__hit_table_rule_outer_tl } }
    \__hit_listing_caption_line:nn
      { \bool_if:NTF \g__hit_en_bool { en } { zh } } {#2}
    \tl_if_empty:NF \l__hit_listing_en_tl
      {
        \bool_if:NT \g__hit_bilingual_caption_bool
          {
            \group_begin:
              \bool_if:NF \g__hit_en_bool {\hit@caption@latin@pick:}
              \exp_args:Ne \__hit_listing_caption_line:nV
                {\bool_if:NTF\g__hit_en_bool{zh}{en}} \l__hit_listing_en_tl
            \group_end:
          }
      }
    \str_case:Vn \g__hit_listing_frame_tl
      {
        { tabular } { \__hit_listing_rule:n { \g__hit_table_rule_outer_tl } }
        { ruled } { \__hit_listing_rule:n { \g__hit_table_rule_inner_tl } }
        { boxed } { \__hit_float_blank:n { table-below } }
      }
  }
\tl_new:N \l__hit_listing_caption_tl
\cs_new_protected:Npn \__hit_listing_caption_line:nn #1#2
  {
    \bool_if:NTF \l__hit_listing_continued_bool
      {
        \tl_set:Nn \l__hit_listing_caption_tl
          { { caption-table } { paragraph = { alignment = right } } }
      }
      {
        \str_if_eq:VnTF \g__hit_listing_frame_tl { ruled }
          {
            \tl_set:Ne \l__hit_listing_caption_tl
              {
                { table-cell }
                { paragraph = { alignment = left , indentation = { left = \dim_use:N \tabcolsep } } }
              }
          }
          { \tl_set:Nn \l__hit_listing_caption_tl { { caption-table } { } } }
      }
    \exp_last_unbraced:NV \docxlike_format_par:nnn \l__hit_listing_caption_tl
      {
        \__hit_listing_label:n {#1}
        \bool_if:NTF \l__hit_listing_continued_bool
          { \hit@term@caption@listing@continued }
          { \hit@caption@separator \ignorespaces #2 }
      }
    \nointerlineskip
  }
\cs_generate_variant:Nn \__hit_listing_caption_line:nn { nV }
\cs_new:Npn \__hit_listing_label:n #1
  {
    \str_if_eq:eeTF {#1} { en }
      {
        \bool_if:NTF \g__hit_en_bool
          { \hit@term@caption@listing@en \nobreakspace }
          { \hit@term@caption@listing@prefix@en }
      }
      { \hit@term@caption@listing@zh }
    \thelstlisting
  }
%</shared-listing>
%    \end{macrocode}
%
% \Finale
